% ===================================================================
%  TABLEAU N° 5 — La Maison et sa construction.
%  Feuillets 31 a 36 du fac-simile = folios 29 a 34.
%  Correspondance : folio imprime = numero de feuillet - 2.
%
%  Comme dans texto/io/14-tabelo-05.tex, ce tableau est un DIALOGUE
%  (Jean, son oncle, M. Leblanc l'entrepreneur, M. Roux l'architecte)
%  ou nul alinea n'est numerote par l'auteur. On klavize donc la
%  totalite du tableau sous le numero d'alinea « 01 », chaque replique
%  ou alinea suivant constituant un rang successif.
%
%  Le francais est plus continu que l'ido : il ne porte, dans les
%  feuillets 31 a 36, qu'un seul intertitre (« Disposition interieure
%  de la maison. ») la ou l'ido en compte sept (« Bona renkontro »,
%  « La laboristi », « La materialo », etc.) — Rochelle laisse courir
%  le texte que Guignon a decoupe en sections (cf. la remarque de
%  texto/fr/10-tableau-01.tex).
%
%  ATTENTION — LE FAC-SIMILE FRANCAIS EST UNE PRISE DE VUE : contraste
%  et echelle inegaux d'une page a l'autre. Chaque nombre de renvoi a
%  ete verifie, quand le fac-simile francais restait ambigu, par
%  comparaison avec le meme passage dans texto/io/14-tabelo-05.tex
%  (numeros de renvoi partages par les deux editions).
% ===================================================================

% --- Feuillet 31 (folio 29, ouverture de tableau) -------------------
\begin{VUpage}[31]{29}
%%K t05-apar-1 apar
\VUsaut{6.0mm}
\VUtitre{15.0pt}{60}{TABLEAU N\textsuperscript{o} 5}
\VUsaut{1.4mm}
\VUtitre{11.4pt}{0}{La Maison et sa construction.}
\VUsaut{2.2mm}

%%K t05-01-1 p
\VUblancAlinea
\textsc{Jean}. --- Mon oncle, je voudrais bien savoir com\cc
ment on bâtit une \VUgras{maison}.

%%K t05-01-2 p
\VUblancAlinea
\textsc{L'Oncle} \textsuperscript{(80)}. --- C'est très facile, mon ami; viens\nl
avec moi, je te montrerai celle que M. Bernard \textsuperscript{(79)}\nl
fait construire et que je dois habiter.

%%K t05-01-3 p
\VUblancAlinea
\textsc{Jean}. --- Qui donc a dessiné le \VUgras{plan} \textsuperscript{(76)} de cette\nl
maison?

%%K t05-01-4 p
\VUblancAlinea
\textsc{L'Oncle}. --- C'est l'\VUgras{architecte} \textsuperscript{(75)}; il a présenté\nl
un dessin très clair, qui a été approuvé, et l'\VUgras{entre}\cc
\VUgras{preneur} \textsuperscript{(77)} a commencé les travaux.

%%K t05-01-5 p
\VUblancAlinea
\textsc{Jean}. --- Qui donc est l'entrepreneur?

%%K t05-01-6 p
\VUblancAlinea
\textsc{L'Oncle}. --- C'est M. Leblanc, le père de ton ami\nl
Jacques. Il dirige les ouvriers avec M. Roux, l'archi\cc
tecte.

%%K t05-01-7 p
\VUblancAlinea
Tiens, voici justement M. Leblanc avec sa \VUgras{règle} \textsuperscript{(78)},\nl
et M. Roux avec son plan.

%%K t05-01-8 p
\VUblancAlinea
Bonjour, Messieurs. Comment vous portez-vous?

%%K t05-01-9 p
\VUblancAlinea
\textsc{M. Roux}. --- Très bien, merci. Bonjour Jean;\nl
comme il a grandi cet enfant!

%%K t05-01-10 p
\VUblancAlinea
\textsc{L'Oncle}. --- Mais oui, c'est un petit homme. Il est\nl
très sérieux; il faut le renseigner sur tout ce qu'il\nl
voit. Aujourd'hui, il voudrait savoir comment on\nl
bâtit une maison.

%%K t05-01-11 p
\VUblancAlinea
\textsc{M. Leblanc}. --- Alors, viens avec moi, mon petit\nl
ami, je vais t'expliquer cela, car j'aime beaucoup les\nl
enfants qui veulent s'instruire.

%%K t05-01-12 p
\VUblancAlinea
Tu vois cet ouvrier qui a une \VUgras{pelle} \textsuperscript{(82)} à la main :\nl
c'est un \VUgras{terrassier} \textsuperscript{(81)}. Il creuse une \VUgras{tranchée} \textsuperscript{(46)}\nl
pour placer la conduite d'eau. C'est lui aussi qui a
\parplein
\end{VUpage}

% --- Feuillet 32 (folio 30) ------------------------------------------
\begin{VUpage}[32]{30}
%%K t05-01-12 p suite
\VUcontinue
creusé le sol à l'endroit où se trouve la maison, pour\nl
que le maçon y élève les quatre \VUgras{murs}. La partie de\nl
ces murs qui est dans la terre s'appelle les \VUgras{fonda}\cc
\VUgras{tions} \textsuperscript{(2)}. L'espace compris entre les fondations forme\nl
la \VUgras{cave} \textsuperscript{(3)}. Cette cave a une petite fenêtre appelée\nl
\VUgras{soupirail} \textsuperscript{(5)}, et une \VUgras{porte} \textsuperscript{(4)}.

%%K t05-01-13 p
\VUblancAlinea
Le \VUgras{maçon} \textsuperscript{(108)} a continué à élever ses murs en\nl
laissant des ouvertures pour les \VUgras{portes} \textsuperscript{(8)} et les\nl
\VUgras{fenêtres} \textsuperscript{(17)}.

%%K t05-01-14 p
\VUblancAlinea
\textsc{Jean}. --- Mais je ne le vois pas, ce maçon!

%%K t05-01-15 p
\VUblancAlinea
\textsc{M. Leblanc}. --- Maintenant il a fini de travailler\nl
à la maison. Il est près de l'\VUgras{écurie} \textsuperscript{(54)}, sur son \VUgras{écha}\cc
\VUgras{faudage} \textsuperscript{(110)}; il construit le mur de la \VUgras{buanderie} \textsuperscript{(56)}.

%%K t05-01-16 p
\VUblancAlinea
Il a une truelle, une auge et \VUgras{un fil à plomb} \textsuperscript{(109)}.\nl
Mais tu ne retiendras pas ces noms-là.

%%K t05-01-17 p
\VUblancAlinea
\textsc{Jean}. --- Si, car je vais demander à mon oncle une\nl
petite truelle et une auge et je ferai moi-même un fil\nl
à plomb avec de la ficelle.

%%K t05-01-18 p
\VUblancAlinea
\textsc{L'Oncle}. --- Ah! très bien. Mais dis-moi, Jean,\nl
que faut-il pour construire une maison?

%%K t05-01-19 p
\VUblancAlinea
\textsc{Jean}. --- Il faut des \VUgras{pierres} \textsuperscript{(59)} de taille et du\nl
\VUgras{mortier} \textsuperscript{(65)}.

%%K t05-01-20 p
\VUblancAlinea
\textsc{L'Oncle}. --- Parfaitement; vois cet homme, avec\nl
son grand chapeau, sur cette \VUgras{charrette} \textsuperscript{(136)}. C'est\nl
un \VUgras{charretier} \textsuperscript{(135)}. Chaque jour il apporte ici des\nl
\VUgras{blocs} \textsuperscript{(60)} de pierre. Il décharge sa voiture dans\nl
la cour, et les \VUgras{tailleurs de pierre} \textsuperscript{(83)} taillent les\nl
blocs et les scient avec leur \VUgras{scie à pierre} \textsuperscript{(85)}. Un\nl
\VUgras{manœuvre} \textsuperscript{(140)} les porte au maçon qui les pose.

%%K t05-01-21 p
\VUblancAlinea
Pendant ce temps, le \VUgras{gâche-mortier} \textsuperscript{(87)} prépare\nl
le mortier.

%%K t05-01-21b p
\VUblancAlinea
Il a besoin de \VUgras{sable} \textsuperscript{(62)}, de \VUgras{chaux} \textsuperscript{(63)} et d'\VUgras{eau} \textsuperscript{(64)}.\nl
La chaux est près de lui dans un \VUgras{sac} \textsuperscript{(71)}; un \VUgras{aide}\cc
\VUgras{maçon} \textsuperscript{(111)} lui apporte le sable dans une \VUgras{brouette} \textsuperscript{(112)}\nl
et l'\VUgras{apprenti} \textsuperscript{(113)} est à la \VUgras{pompe} \textsuperscript{(58)}. Il remplit un\nl
\VUgras{seau} \textsuperscript{(114)}. Le mortier est bientôt prêt. Le \VUgras{porte}\cc
\parplein
\end{VUpage}

% --- Feuillet 33 (folio 31) ------------------------------------------
\begin{VUpage}[33]{31}
%%K t05-01-21b p suite
\VUcontinue
\VUgras{mortier} \textsuperscript{(107)} le prend et le monte par l'échelle sur\nl
l'échafaudage, où travaille un maçon qui termine ce\nl
côté de la maison.

%%K t05-01-22 p
\VUblancAlinea
Et maintenant, comment s'appelle l'ouvrier qui fait\nl
le \VUgras{toit} \textsuperscript{(31)} ?

%%K t05-01-23 p
\VUblancAlinea
\textsc{Jean}. --- C'est le \VUgras{couvreur} \textsuperscript{(118)}.

%%K t05-01-24 p
\VUblancAlinea
\textsc{M. Leblanc}. --- Oui, mais il y a d'abord le \VUgras{char}\cc
\VUgras{pentier} \textsuperscript{(115)}.

%%K t05-01-25 p
\VUblancAlinea
Le charpentier fait la \VUgras{charpente} \textsuperscript{(35)}. Il assemble\nl
les \VUgras{poutres} \textsuperscript{(37)} avec des \VUgras{chevilles} \textsuperscript{(117)}. Ces ouvriers,\nl
là-haut, avec de larges culottes de velours, sont des\nl
charpentiers. L'un tient une petite poutre et l'autre\nl
coupe une cheville avec sa \VUgras{hache} \textsuperscript{(116)}. Celui qui est\nl
près de la \VUgras{cheminée} \textsuperscript{(41)} est le couvreur.

%%K t05-01-25b p
\VUblancAlinea
Il cloue des lattes sur les solives et des \VUgras{ardoi}\cc
\VUgras{ses} \textsuperscript{(66)} sur les lattes. Quelquefois il pose des tuiles.

%%K t05-01-26 p
\VUblancAlinea
Le toit de l'écurie est en \VUgras{tuiles} \textsuperscript{(55)}, celui de la\nl
maison est en \VUgras{ardoises} \textsuperscript{(32)} et celui de la véranda est\nl
en \VUgras{zinc} \textsuperscript{(24)}.

%%K t05-01-27 p
\VUblancAlinea
\textsc{Jean}. --- Est-ce que le couvreur fait aussi les toits\nl
en zinc ?

%%K t05-01-28 p
\VUblancAlinea
\textsc{M. Leblanc}. --- Non, c'est le \VUgras{zingueur} \textsuperscript{(121)} ou le\nl
\VUgras{plombier}, cet ouvrier qui pose une \VUgras{lucarne} \textsuperscript{(38)} sur\nl
le toit de la maison. Il pose aussi les \VUgras{chéneaux} \textsuperscript{(43)} et\nl
les \VUgras{gouttières} \textsuperscript{(42)}. Il soude le zinc et le fer-blanc.\nl
C'est lui qui a fixé le \VUgras{paratonnerre} \textsuperscript{(40)} et la\nl
\VUgras{girouette} \textsuperscript{(39)} sur le \VUgras{pignon} \textsuperscript{(36)}.

%%K t05-01-29 p
\VUblancAlinea
\textsc{Jean}. --- Alors la maison est finie?

%%K t05-01-30 p
\VUblancAlinea
\textsc{M. Leblanc}. --- Pas tout à fait. Le \VUgras{plâtrier} \textsuperscript{(99)}\nl
construit avec des \VUgras{briques} \textsuperscript{(61)} les cloisons qui la divi\cc
sent en différentes pièces. Il enduit de \VUgras{plâtre} \textsuperscript{(70)}\nl
les \VUgras{plafonds} \textsuperscript{(26)} et les murailles. Il fait en ce moment\nl
le plafond de la \VUgras{véranda} \textsuperscript{(33)}. Il a, comme le maçon,\nl
une \VUgras{truelle} \textsuperscript{(100)} et une \VUgras{auge} \textsuperscript{(101)}.

%%K t05-01-31 p
\VUblancAlinea
Puis le menuisier fait les portes, les fenêtres, les\nl
\VUgras{parquets} \textsuperscript{(16)} et les \VUgras{volets} \textsuperscript{(19)}.
\end{VUpage}

% --- Feuillet 34 (folio 32) ------------------------------------------
\begin{VUpage}[34]{32}
%%K t05-01-32 p
\VUblancAlinea
Il travaille maintenant dans la cour sur son\nl
\VUgras{établi} \textsuperscript{(90)}. Il a une \VUgras{scie} \textsuperscript{(94)} pour scier le \VUgras{bois} \textsuperscript{(69)}, un\nl
\VUgras{rabot} \textsuperscript{(91)}, un \VUgras{valet} \textsuperscript{(92)}, un \VUgras{ciseau} \textsuperscript{(94 \textit{bis})}, un\nl
\VUgras{compas} \textsuperscript{(95)} et un \VUgras{maillet} \textsuperscript{(95 \textit{bis})}.

%%K t05-01-33 p
\VUblancAlinea
\textsc{Jean}. --- Ah! mais c'est encore plus amusant d'être\nl
menuisier que d'être maçon. Mon oncle, tu m'achè\cc
teras un établi, une scie et un rabot, car je vais faire\nl
une niche pour Médor.

%%K t05-01-34 p
\VUblancAlinea
\textsc{L'Oncle}. --- Oui, mon enfant.

%%K t05-01-35 p
\VUblancAlinea
\textsc{Jean}. --- Que je suis content! Et celui-là, près de\nl
la porte d'entrée, qui est-ce?

%%K t05-01-36 p
\VUblancAlinea
\textsc{M. Leblanc}. --- C'est un \VUgras{peintre} \textsuperscript{(102)}. Il est sur\nl
son échelle avec son pot de \VUgras{peinture} \textsuperscript{(103)} et son\nl
\VUgras{pinceau} \textsuperscript{(104)}. Il enduit le bois de la porte d'entrée\nl
pour le conserver plus longtemps.

%%K t05-01-37 p
\VUblancAlinea
C'est lui aussi qui colle les papiers dans les appar\cc
tements.

%%K t05-01-38 p
\VUblancAlinea
Le \VUgras{tapissier} \textsuperscript{(107)} qui est sur une \VUgras{échelle} \textsuperscript{(106)},\nl
sur le \VUgras{balcon} \textsuperscript{(28)}, pose un \VUgras{store} \textsuperscript{(29)}.

%%K t05-01-39 p
\VUblancAlinea
L'ouvrier près de la grande fenêtre est le \VUgras{vitrier} \textsuperscript{(96)}.\nl
Il pose les \VUgras{vitres} \textsuperscript{(18)} avec du \VUgras{mastic} \textsuperscript{(73)}. Cet autre\nl
ouvrier, agenouillé près de la porte de la cave, est un\nl
\VUgras{serrurier} \textsuperscript{(97)}. Il place une \VUgras{serrure} \textsuperscript{(6)}. Sa \VUgras{lime} \textsuperscript{(98)}\nl
est près de lui.

%%K t05-01-40 p
\VUblancAlinea
\textsc{Jean}. --- Est-ce aussi un serrurier, celui qui frappe\nl
avec son \VUgras{marteau} \textsuperscript{(84)} sur un morceau de fer?

%%K t05-01-41 p
\VUblancAlinea
\textsc{M. Leblanc}. --- C'est plutôt un \VUgras{forgeron} \textsuperscript{(125)}. Il\nl
forge des \VUgras{morceaux de fer} \textsuperscript{(67)} pour \VUgras{l'électri}\cc
\VUgras{cien} \textsuperscript{(124)} qui est sur le toit de la \VUgras{remise} \textsuperscript{(51)}. Il a une\nl
\VUgras{forge} \textsuperscript{(126)}, une \VUgras{enclume} \textsuperscript{(127)} et une \VUgras{pince} \textsuperscript{(128)}.

%%K t05-01-42 p
\VUblancAlinea
L'électricien pose le \VUgras{fil de cuivre} \textsuperscript{(44)} du téléphone.

%%K t05-01-43 p
\VUblancAlinea
\textsc{L'Oncle}. --- Au fond de la \VUgras{cour} \textsuperscript{(45)}, près de la\nl
\VUgras{grille} \textsuperscript{(47)} et de la \VUgras{porte cochère} \textsuperscript{(48)}, tu vois le \VUgras{jar}\cc
\VUgras{dinier} \textsuperscript{(129)}. Il prépare le petit \VUgras{jardin} \textsuperscript{(49)}. Il a bêché\nl
la terre avec sa \VUgras{bêche} \textsuperscript{(131)}; il sème des graines et\nl
ratisse avec le \VUgras{râteau} \textsuperscript{(130)}. Puis, avec son \VUgras{arro}\cc
\parplein
\end{VUpage}

% --- Feuillet 35 (folio 33) ------------------------------------------
\begin{VUpage}[35]{33}
%%K t05-01-43 p suite
\VUcontinue
\VUgras{soir} \textsuperscript{(132)}, il arrosera les \VUgras{plates-bandes} \textsuperscript{(150)} et les\nl
plantes pousseront.

%%K t05-01-44 p
\VUblancAlinea
Le \VUgras{valet d'écurie} \textsuperscript{(133)}, qui vient de soigner le\nl
cheval et de lui donner à manger, nettoie la \VUgras{voiture} \textsuperscript{(52)}\nl
avec une éponge, une \VUgras{pompe à main} \textsuperscript{(134)} et un seau\nl
d'eau. Puis il la rentrera dans la remise et fermera\nl
la porte avec le gros \VUgras{verrou} \textsuperscript{(53)}.

%%K t05-01-45 p
\VUblancAlinea
Te voilà satisfait, j'espère, tu as eu assez d'explica\cc
tions !

%%K t05-01-46 p
\VUblancAlinea
\textsc{Jean}. --- Oui, mon oncle; mais je voudrais voir\nl
l'intérieur de la maison.

%%K t05-01-47 p
\VUblancAlinea
\textsc{L'Oncle}. --- Ce sera pour demain. M. Roux t'expli\cc
quera le plan et t'indiquera le nom de chaque pièce.

%%K t05-tit-1 sub
\VUsaut{3.0mm}
\VUpk{10.2pt}{40}{Disposition intérieure de la maison.}
\VUsaut{2.4mm}

%%K t05-apar-2 apar
{\centering\textit{(Voir le Plan.)}\par}
\VUsaut{2.4mm}

%%K t05-01-48 p
\VUblancAlinea
\textsc{L'Architecte}. --- Viens, Jean, je vais te faire\nl
visiter les différentes parties de la maison.

%%K t05-01-49 p
\VUblancAlinea
Entrons au \VUgras{rez-de-chaussée} \textsuperscript{(7)}. Voilà le bouton\nl
de la \VUgras{sonnette} \textsuperscript{(11)}. Nous montons les trois \VUgras{mar}\cc
\VUgras{ches} \textsuperscript{(14)} du \VUgras{perron} \textsuperscript{(9)}, nous passons le \VUgras{seuil} \textsuperscript{(10)} de\nl
la \VUgras{porte d'entrée} \textsuperscript{(8)} et nous voici au pied de l'\VUgras{esca}\cc
\VUgras{lier} \textsuperscript{(13)}.

%%K t05-01-50 p
\VUblancAlinea
\textsc{Jean}. --- C'est le corridor.

%%K t05-01-51 p
\VUblancAlinea
\textsc{L'Architecte}. --- Non, c'est le \VUgras{vestibule} \textsuperscript{(12)}. Le\nl
\VUgras{corridor} \textsuperscript{(a)} se trouve au bout. A droite, voici le\nl
\VUgras{salon} \textsuperscript{(b)} et la \VUgras{salle à manger} \textsuperscript{(c)}; en face de la\nl
salle à manger se trouvent la \VUgras{cuisine} \textsuperscript{(d)} et l'\VUgras{office} \textsuperscript{(e)};\nl
puis sur le devant, le \VUgras{cabinet de travail} \textsuperscript{(f)}. La\nl
\VUgras{véranda} \textsuperscript{(23)}, avec sa \VUgras{porte vitrée} \textsuperscript{(25)}, est auprès du\nl
salon.

%%K t05-01-52 p
\VUblancAlinea
Nous allons monter l'escalier. Tiens la \VUgras{rampe} \textsuperscript{(15)}\nl
et compte les marches. Il y en a quatorze. Voici le\nl
\VUgras{palier} \textsuperscript{(m)}, nous sommes au \VUgras{premier étage} \textsuperscript{(27)}.
\end{VUpage}

% --- Feuillet 36 (folio 34) ------------------------------------------
\begin{VUpage}[36]{34}
%%K t05-01-53 p
\VUblancAlinea
En face de l'escalier, il y a les \VUgras{cabinets d'ai}\cc
\VUgras{sances} \textsuperscript{(20)} et le \VUgras{cabinet de toilette} \textsuperscript{(j)}. A droite,\nl
une belle \VUgras{chambre à coucher} \textsuperscript{(g)} avec deux fenêtres\nl
sur la \VUgras{façade} \textsuperscript{(1)} de la maison. A gauche, c'est la \VUgras{salle}\nl
\VUgras{de bains} \textsuperscript{(1)} et une \VUgras{chambre d'amis} \textsuperscript{(i)} avec une\nl
grande \VUgras{alcôve} \textsuperscript{(h)}. A côté de la chambre à coucher, il\nl
y a la \VUgras{chambre des enfants} \textsuperscript{(k)}. Elle donne sur\nl
une vaste \VUgras{terrasse} \textsuperscript{(21)}. Sous cette terrasse sont\nl
d'autres \VUgras{cabinets d'aisances} \textsuperscript{(20)} et, contre le mur,\nl
voici la \VUgras{cloche} \textsuperscript{(22)} pour sonner le dîner.

%%K t05-01-54 p
\VUblancAlinea
\textsc{Jean}. --- Allons au \VUgras{deuxième étage}.

%%K t05-01-55 p
\VUblancAlinea
\textsc{L'Architecte}. --- Il n'y a pas de deuxième étage.\nl
Sous le toit, il n'y a que des \VUgras{mansardes} \textsuperscript{(33)} et le\nl
\VUgras{grenier} \textsuperscript{(34)}. Nous avons fini la visite, nous pouvons\nl
redescendre.

%%K t05-01-56 p
\VUblancAlinea
\textsc{Jean}. --- Merci, Monsieur, c'est très intéressant.\nl
Je vais raconter à papa tout ce que j'ai vu.

\VUsaut{4.0mm}
\VUfilet{22mm}
\end{VUpage}
